% Generated draft: requires mathematical and source review.
\documentclass[11pt]{article}
\usepackage[margin=25mm]{geometry}
\usepackage{fontspec}
\setmainfont{Latin Modern Roman}
\usepackage{amsmath,amssymb,mathtools,unicode-math}
\setmathfont{Latin Modern Math}
\usepackage{enumitem,needspace,xurl,hyperref}
\setlength{\parindent}{0pt}
\setlength{\parskip}{5pt}
\setlength{\emergencystretch}{3em}
\newcommand{\N}{\mathbb{N}}
\newcommand{\Z}{\mathbb{Z}}
\newcommand{\Q}{\mathbb{Q}}
\newcommand{\R}{\mathbb{R}}
\newcommand{\C}{\mathbb{C}}
\newcommand{\F}{\mathbb{F}}
\newcommand{\A}{\mathbb{A}}
\newcommand{\PP}{\mathbb{P}}
\newcommand{\E}{\mathbb{E}}
\newcommand{\eps}{\varepsilon}
\newcommand{\dd}{\,\mathrm{d}}
\begin{document}
\section*{UnsolvedMath ID: 30005571}
\textbf{Source problem code:} catalogue entry
\subsection{Definitions and mathematical statement}
The supplied record concerns the cyclotomic \(K\)-theory modular-symbol maps denoted by \(\widetilde\Pi\) and \(\Pi\), and an Eisenstein ideal \(I\) in the relevant Hecke algebra. The record does not specify that algebra, the generators of \(I\), the domains and codomains of the maps, or their coefficient rings. These names therefore refer to the source construction; they do not define a complete construction here.

For an additive map \(f\colon M\to A\) whose domain carries an action of an ideal \(I\), the Eisenstein property in the sense explicitly used in the source is
\[
 f\circ T=0\quad\text{for every }T\in I,
 \qquad\text{equivalently}\qquad
 f(Tx)=0\quad\text{for every }T\in I\text{ and }x\in M.
\]
Here \(M\) and \(A\) are explanatory placeholders, not identified source modules. Writing \(IM\) for the additive subgroup generated by the elements \(Tx\), this property is equivalent to \(IM\subseteq\ker(f)\), and gives an induced additive map \(M/IM\to A\).

The source's two requested assertions are:
\begin{enumerate}
\item The map \(\widetilde\Pi\) is Eisenstein:
\[
 \widetilde\Pi\circ T=0\qquad\text{for every }T\in I.
\]
\item The map denoted in the source by \(\Pi\bmod I\) is an isomorphism.
\end{enumerate}
The second assertion cannot be made into a fully specified isomorphism statement from this record: the source does not identify the modules, quotient construction, or category intended by \(\Pi\bmod I\). In particular, the first assertion concerns \(\widetilde\Pi\), whereas the second concerns \(\Pi\); their relationship is not supplied.

The excerpt also reports partial results, retained here as attributed source claims rather than independently verified theorems. It attributes to Fukaya--Kato an earlier form of the conjecture and progress toward the isomorphism assertion, referring to a 2012 preprint [2] and to [5]. Its displayed Theorem 1 asserts that
\[
 \widetilde\Pi\otimes\mathbb Z_p\text{ is Eisenstein for }p\mid N.
\]
The divisibility condition \(p\mid N\) is retained exactly as supplied; the surrounding sentence is corrupted and the tensor-product base is unspecified. The excerpt attributes to S.--Venkatesh the assertion
\[
 \Pi\circ\bigl(T_\ell-\ell-\langle\ell\rangle\bigr)=0
 \qquad\text{for primes }\ell\nmid Np.
\]
The auxiliary prime \(p\), the Hecke and diamond operators, and the coefficient conventions are not defined in the supplied passage.

The supplied literature assessment, dated 2026-08-22, labels the problem partially solved: it reports Eisenstein results in substantial ranges, including away from the level and, after inverting \(3\), for the cyclotomic map, while describing the asserted isomorphism modulo the Eisenstein ideal as conjectural. The record does not specify the exact hypotheses or identify which map is meant in every part of that assessment. This is supplied dated curation, not an independently verified mathematical or literature assessment.
\subsection{Short English statement}
The source asks whether every operator in the Eisenstein ideal annihilates the map \(\widetilde\Pi\), and whether the map written \(\Pi\bmod I\) is an isomorphism. Exact map definitions, coefficients, level restrictions, and the meaning of the second quotient are missing. The supplied dated assessment reports partial results and leaves the isomorphism assertion conjectural; this draft requires mathematical review.
\subsection{Sources}
\begin{itemize}
\item Supplied local UnsolvedMath registry record, permanent ID \(30005571\), source code \texttt{catalogue entry}; registry collection \texttt{data/unsolvedmath/problems.json}, with \texttt{catalog/problems.json} designated for additional existing records. This draft uses the supplied record and does not claim to have inspected those files.
\item \emph{Algebraische Zahlentheorie} (2023), Oberwolfach Reports, pp.\ 1543--1606, DOI 10.4171/owr/2023/28. Supplied report reference: \url{https://doi.org/10.4171/owr/2023/28}. No exact page or speaker attribution for this conjecture is supplied.
\item Sharifi--Venkatesh, \emph{Eisenstein cocycles in motivic cohomology}: \url{https://doi.org/10.1112/S0010437X24007322}. Supplied literature-review reference.
\item Lecouturier--Sharifi--Shih--Wang, \emph{Eisenstein cocycles for imaginary quadratic fields}: \url{https://arxiv.org/abs/2504.19125}. Supplied literature-review reference; its precise bearing on the cyclotomic assertions is not established by the record.
\item The excerpt attributes partial results to Fukaya--Kato and S.--Venkatesh. It refers to a 2012 Fukaya--Kato preprint [2], a paper by Fukaya, Kato, and the unnamed speaker [3], earlier work of Goncharov [4], and references [5] and [6]. Complete bibliographic entries and a verified correspondence with the supplied literature links are absent.
\item Literature assessment checked at \texttt{2026-08-22}: the label \texttt{partially\_solved} and its accompanying claims are supplied dated curation, not independently verified here. The registry also supplies the label \texttt{corrected\_verified} and says its clean statement was checked against the official report; that verification claim is retained as registry metadata, not adopted as verification performed for this draft.
\end{itemize}
\end{document}
